% SST Project Report - Section 6 deliverable draft for Overleaf.
% Replace bracketed placeholders and the screenshot frame before submission.
% Suggested compiler: pdfLaTeX. Target length: approximately 2-4 pages.

\documentclass[11pt,a4paper]{article}
\usepackage[margin=2cm]{geometry}
\usepackage[T1]{fontenc}
\usepackage[utf8]{inputenc}
\usepackage{lmodern}
\usepackage{graphicx}
\usepackage{tabularx}
\usepackage{array}
\usepackage{enumitem}
\usepackage{xcolor}
\usepackage{hyperref}
\usepackage{titlesec}

\hypersetup{
  colorlinks=true,
  linkcolor=blue,
  urlcolor=blue,
  pdftitle={Student Study Tracker Project Report},
  pdfauthor={[Student Name]}
}
\setlist{nosep,leftmargin=*}
\setlength{\parskip}{4pt}
\setlength{\parindent}{0pt}
\titleformat{\section}{\large\bfseries\color{blue!55!black}}{\thesection}{0.6em}{}
\titleformat{\subsection}{\normalsize\bfseries}{\thesubsection}{0.6em}{}

\newcommand{\placeholder}[1]{\textcolor{gray}{\textit{[#1]}}}

\begin{document}

\begin{center}
  {\LARGE\bfseries Section 6: Student Study Tracker Project Deliverable}\\[5pt]
  {\large Student Study Tracker (SST)}\\[9pt]
  \begin{tabular}{rl}
    Student: & \placeholder{Student name} \\
    Programme / course: & \placeholder{Programme and course} \\
    Institution: & Accra Technical University \\
    Submission date: & \placeholder{Final submission date}
  \end{tabular}
\end{center}

\section*{Project Links}
\textbf{Live application URL:} \url{https://sstv1-e6dc0.web.app}\\
\textbf{GitHub repository URL:} \url{https://github.com/Gideon-Gyebi/sst-internship-project}

\section{Project Overview}

\textbf{Problem statement.} Students may keep coursework and study actions in separate notes or reminders, making it difficult to see what remains to be done and what has been completed. The project addresses this problem with a simple, secure study-task tracker that organises work by title and course.

\textbf{Plan.} Build a responsive dashboard where a student can add tasks with a title and course, edit or delete them, mark them complete, and review progress. Use Firebase Authentication to separate student accounts, Cloud Firestore to store each student's tasks, and Firestore Security Rules to enforce ownership and validate task data. The project was developed for the ATU Virtual Internship beginner software development project.

\section{Implemented Features}

SST uses Vite, HTML, CSS, and modular vanilla JavaScript for its client interface. Firebase Authentication supports email and password accounts and Google sign-in. Cloud Firestore stores task data under a user-specific path. Firebase Hosting serves the deployed web application.

The current application supports the following:
\begin{itemize}
  \item Account registration, sign-in, restored authentication state, and sign-out.
  \item Task creation and display with a required title and subject/course.
  \item Editing, deletion, completion, and reopening of a task.
  \item Dashboard counts for total, pending, and completed tasks.
  \item All, Pending, and Completed status views, with a maximum of ten tasks shown per page.
  \item Loading, retry, error, and empty-state feedback.
  \item Responsive layouts and a saved manual light/dark theme preference.
\end{itemize}

Each task contains a title, subject, status, creation timestamp, and update timestamp. Firestore Security Rules restrict task access to the authenticated owner and validate allowed fields, value types, status values, size limits, and server timestamps. Whitespace-only titles and course names are rejected by the rules.

\section{Feature Screenshot Evidence}

\begin{figure}[htbp]
  \centering
  \fbox{\begin{minipage}[c][4.2cm][c]{0.88\linewidth}
    \centering
    \textbf{Dashboard screenshot placeholder}\\[5pt]
    Insert a dashboard screenshot showing the core features, including at least one added task and one completed task.
  \end{minipage}}
  \caption{Implemented SST features displayed on the dashboard. Replace the placeholder with a clear application screenshot.}
  \label{fig:sst-dashboard}
\end{figure}

\section{Deployment and Evidence}

The Firestore rules and Hosting build were released to the Firebase project for SST on 31 October 2026. The deployed sign-in page was opened successfully after release. Authenticated task operations have not yet been retested against the live deployment.

\section{Testing Checklist}

The Firestore Emulator Suite ran seven automated security-rule tests using an isolated demo project. The production build, dependency audit, pagination, and initial live page load were also checked; authenticated task operations still need a live retest. The checklist records the available results and evidence still to be confirmed.

\begin{center}
\small
\begin{tabularx}{\linewidth}{>{\raggedright\arraybackslash}p{0.40\linewidth} X}
\hline
\textbf{Test} & \textbf{Result / evidence} \\
\hline
Firestore access and validation rules & Pass: 7 emulator tests; no production writes \\
Production build & Pass: Vite generated the Hosting build \\
Dependency audit & Pass: zero reported vulnerabilities \\
Pagination & Pass: ten-item page limit observed locally \\
Live deployment & Pass: deployed sign-in page loaded; authenticated operations pending retest \\
Responsive layout and light/dark themes & \placeholder{Owner to confirm device and theme checks} \\
\hline
\end{tabularx}
\end{center}

\begin{figure}[htbp]
  \centering
  \fbox{\begin{minipage}[c][3.5cm][c]{0.46\linewidth}
    \centering
    \textbf{VS Code testing screenshot}\\[5pt]
    Show the project codebase and test/build results in the terminal.
  \end{minipage}}\hfill
  \fbox{\begin{minipage}[c][3.5cm][c]{0.46\linewidth}
    \centering
    \textbf{GitHub screenshot}\\[5pt]
    Show the project repository and relevant verification evidence.
  \end{minipage}}
  \caption{Testing and repository evidence. Replace both placeholders with screenshots.}
  \label{fig:testing-evidence}
\end{figure}

\section{Development Challenge and Resolution}

A key challenge was protecting personal task data while keeping the interface responsive to Firestore updates. Client-side checks alone could be bypassed by a modified client, so the project enforces ownership and task validation in Firestore Security Rules as well. Emulator tests now exercise both allowed and denied requests without touching live data. During dashboard testing, pagination was also checked for keyboard focus retention after each page change.

\section{Future Improvement}

A future improvement is a combined search and course filter. Students could search task titles or descriptions and narrow results to a selected course. This feature is outside the current implementation scope; pagination keeps long task lists manageable in the delivered version.

\section*{Items to Complete Before Submission}

\begin{itemize}
  \item Replace the student, programme/course, and final submission date placeholders.
  \item Replace the dashboard, VS Code, and GitHub screenshot placeholders with clear evidence images.
  \item Confirm the responsive and light/dark theme test row after owner review.
  \item Recheck authenticated task operations on the deployed site if required by the instructor.
  \item Add any required AI-assistance disclosure and source citations under ATU guidance.
  \item Confirm whether the instructor requires any additional report headings or citation style.
\end{itemize}

\end{document}
